\documentclass[11pt]{amsart}

\usepackage[T1]{fontenc}
\usepackage[utf8]{inputenc}
\usepackage{lmodern}
\usepackage{microtype}
\usepackage{amsmath,amssymb,amsthm,mathtools,amscd}
\usepackage{enumitem}
\usepackage[colorlinks=true,linkcolor=blue,citecolor=blue,urlcolor=blue]{hyperref}
\usepackage[margin=1in]{geometry}

\numberwithin{equation}{section}

\newtheorem{theorem}{Theorem}[section]
\newtheorem{proposition}[theorem]{Proposition}
\newtheorem{lemma}[theorem]{Lemma}
\newtheorem{corollary}[theorem]{Corollary}
\theoremstyle{definition}
\newtheorem{definition}[theorem]{Definition}
\newtheorem{axiom}[theorem]{Axiom}
\newtheorem{example}[theorem]{Example}
\theoremstyle{remark}
\newtheorem{remark}[theorem]{Remark}

\newcommand{\Noetica}{\mathrm{Noetica}}
\newcommand{\NoeTop}{\mathrm{NoeTop}}
\newcommand{\FinSet}{\mathrm{FinSet}}
\newcommand{\Sh}{\mathrm{Sh}}
\newcommand{\id}{\mathrm{id}}
\newcommand{\I}{\mathcal{I}}
\newcommand{\IA}{\mathcal{I}_A}
\newcommand{\IB}{\mathcal{I}_B}
\newcommand{\IC}{\mathcal{I}_C}
\newcommand{\ID}{\mathcal{I}_D}
\newcommand{\Jnoe}{J_{\mathrm{noe}}}
\newcommand{\Deltaopen}{\Delta_{39}^{\circ}}
\newcommand{\R}{\mathbb{R}}

\title[The Noetic Topos]{The Noetic Topos: A Categorical Framework for Structured Cognitive Assessment of AI Systems}
\author[Henderson]{Little Henderson}
\date{}

\begin{document}

\begin{abstract}
We propose an axiomatic categorical framework intended as a formal language for structured cognitive assessment of AI systems. Starting from finite distinguished data consisting of forty labeled indices, admissibility relations, and concrete generating maps between finite sets, we define a pullback-complete concrete category \(\Noetica\). We then add two further layers: an ACBE descent structure, encoding designated cross-domain covering families, and an RPM section-retraction pair \(\eta : \id_{\Noetica} \Rightarrow R\), \(\varepsilon : R \Rightarrow \id_{\Noetica}\). These induce a coverage \(\Jnoe\) whose covering sieves are generated by jointly surjective families, designated descent families, or RPM realization maps. The main theorem proves that \(\Jnoe\) is a Grothendieck topology. Consequently, the associated sheaf category \(\Sh(\Noetica,\Jnoe)\) is a Grothendieck topos and hence an elementary topos with exponentials, a subobject classifier, and internal intuitionistic logic. We also exhibit an explicit finite model showing that the axioms are jointly satisfiable, and we attach an information-geometric layer by using the open probability simplex \(\Delta_{39}^{\circ}\) and the Fisher--Rao metric to define noetic manifolds. The paper makes no claim that its axioms are empirically derived, uniquely correct, or already validated for AI evaluation. Its narrower contribution is formal: it isolates a coherent site/topos construction and a compatible geometric layer that could later be tested as bookkeeping machinery for empirical assessment data.
\end{abstract}

\maketitle

\section{Introduction}

AI alignment work increasingly relies on heterogeneous evidence: red-team behavior, task performance, preference data, safety evaluations, interpretability signals, and context-sensitive monitoring. A recurring difficulty is structural rather than numerical. Local observations are made under different prompts, tasks, roles, or deployment conditions, and it is not automatic that they compose coherently.

Category theory is useful in precisely this setting because it distinguishes objects, morphisms, pullbacks, coverings, and gluing conditions. A site keeps track of which local views count as jointly informative, while the associated sheaf category enforces compatibility constraints on data collected over those views. If one wants a formal bookkeeping language for structured AI assessment, a site/topos construction is therefore a natural candidate.

This paper does not claim to derive such a framework from empirical data. Instead, it isolates an axiomatic core and proves its internal mathematical consequences. The contribution is fourfold.
\begin{enumerate}[label=(\roman*)]
\item We define a finite concrete category \(\Noetica\) from explicit distinguished data and pullback closure.
\item We specify ACBE descent data and RPM realization data, the latter as a section-retraction pair rather than an adjunction.
\item We define a coverage \(\Jnoe\) and prove in Theorem~\ref{thm:Jnoe-topology} that it is a Grothendieck topology.
\item We exhibit an explicit finite model showing that the axioms are jointly satisfiable.
\item We attach an information-geometric layer based on the open probability simplex and the Fisher--Rao metric.
\end{enumerate}

The resulting manuscript is deliberately conservative. It proves only formal results: existence of a small site, existence of a sheaf topos, and compatibility of a geometric state-space layer. Any claim that the framework measures cognition, alignment, or safety in a predictive scientific sense is deferred to future empirical work.

\section{Preliminaries}

We recall the standard notions used later; see \cite{MacLane1998,MacLaneMoerdijk1992}.

\begin{definition}
Let \(\mathcal{C}\) be a category and \(C \in \mathrm{Ob}(\mathcal{C})\).
\begin{enumerate}[label=(\roman*)]
\item A \emph{sieve} \(S\) on \(C\) is a collection of morphisms with codomain \(C\) that is closed under precomposition: if \(f : D \to C\) lies in \(S\) and \(h : E \to D\) is any morphism, then \(f \circ h \in S\).
\item If \(g : D \to C\) is a morphism and \(S\) is a sieve on \(C\), the \emph{pullback sieve} \(g^{*}S\) on \(D\) is
\[
g^{*}S := \{h : E \to D \mid g \circ h \in S\}.
\]
\item The \emph{maximal sieve} \(m_C\) on \(C\) is the class of all morphisms with codomain \(C\).
\end{enumerate}
\end{definition}

\begin{definition}
A \emph{Grothendieck topology} \(J\) on a category \(\mathcal{C}\) assigns to each object \(C\) a collection \(J(C)\) of sieves on \(C\), called covering sieves, such that:
\begin{enumerate}[label=(J\arabic*)]
\item \label{J1} the maximal sieve \(m_C\) lies in \(J(C)\);
\item \label{J2} if \(S \in J(C)\) and \(g : D \to C\), then \(g^{*}S \in J(D)\);
\item \label{J3} if \(S \in J(C)\) and \(T\) is a sieve on \(C\) such that \(f^{*}T \in J(D)\) for every \(f : D \to C\) in \(S\), then \(T \in J(C)\).
\end{enumerate}
The pair \((\mathcal{C},J)\) is called a \emph{site}.
\end{definition}

\begin{definition}
Let \((\mathcal{C},J)\) be a site. A presheaf \(F : \mathcal{C}^{\mathrm{op}} \to \mathrm{Set}\) is a \emph{sheaf} if for every object \(C\) and every covering sieve \(S \in J(C)\), the canonical map
\[
F(C) \longrightarrow \mathrm{Nat}(S,F)
\]
is a bijection, where \(S\) is regarded as a subfunctor of the representable presheaf \(yC = \mathrm{Hom}_{\mathcal{C}}(-,C)\).
\end{definition}

\begin{remark}
If \((\mathcal{C},J)\) is a small site, then \(\Sh(\mathcal{C},J)\) is a Grothendieck topos \cite[Chap.~VII]{MacLaneMoerdijk1992}. Consequently it is also an elementary topos, with finite limits, finite colimits, exponentials, a subobject classifier, and intuitionistic internal logic.
\end{remark}

\section{The Category Noetica}

We fix a finite index set \(\I\) with \(|\I|=40\), together with a partition
\[
\I = \IA \sqcup \IB \sqcup \IC \sqcup \ID
\]
and a fixed preorder on the four labels
\[
A \succeq B \succeq C \succeq D.
\]
The letters \(A,B,C,D\) are formal domain labels in this paper. No empirical or metaphysical interpretation is assumed.

\begin{definition}
The distinguished data consist of:
\begin{enumerate}[label=(\roman*)]
\item for each \(i \in \I\), a finite set \(X_i\);
\item for each \(k \in \{0,1,\dots,9\}\), an admissibility relation \(D_k \subseteq \I \times \I\);
\item for each \(k\) and each \((i,j) \in D_k\), a function
\[
\nu_{k,i,j} : X_i \to X_j;
\]
\item the identity condition
\[
D_0 = \{(i,i) : i \in \I\}, \qquad \nu_{0,i,i} = \id_{X_i}.
\]
\end{enumerate}
\end{definition}

\begin{definition}\label{def:noetica}
The category \(\Noetica\) is the smallest subcategory of \(\FinSet\) satisfying the following closure conditions:
\begin{enumerate}[label=(\roman*)]
\item its distinguished objects include every \(X_i\) for \(i \in \I\);
\item its generating morphisms include every \(\nu_{k,i,j}\) with \((i,j)\in D_k\);
\item it is closed under identities and composition;
\item whenever \(f : X \to Z\) and \(g : Y \to Z\) lie in \(\Noetica\), the fiber product \(X \times_Z Y\) is an object of \(\Noetica\), the pullback projections belong to \(\Noetica\), and any morphism obtained from the universal property of this pullback belongs to \(\Noetica\).
\end{enumerate}
Morphisms in \(\Noetica\) are therefore concrete functions between finite sets, and equality of morphisms is extensional.
\end{definition}

\begin{axiom}
The pullback-closure process in Definition~\ref{def:noetica} yields only finitely many objects up to the identifications intended by the construction.
\end{axiom}

\begin{proposition}
Under the preceding axiom, \(\Noetica\) is a small concrete category with pullbacks and finite hom-sets.
\end{proposition}

\begin{proof}
Because every object of \(\Noetica\) is a finite set and every morphism is an actual function, the inclusion
\[
U : \Noetica \hookrightarrow \FinSet
\]
is faithful. Pullbacks exist in \(\Noetica\) by Definition~\ref{def:noetica}(iv). Smallness follows from the axiom that only finitely many objects are produced up to the designated identifications. For objects \(X,Y \in \mathrm{Ob}(\Noetica)\), the hom-set \(\mathrm{Hom}_{\Noetica}(X,Y)\) is a subset of the finite set \(Y^{X}\), hence is finite.
\end{proof}

\section{ACBE Descent Structure}

In this section, ``ACBE'' is only a name for the fixed four-level preorder \(A \succeq B \succeq C \succeq D\).

\begin{axiom}
For every pair of levels \(W',W \in \{A,B,C,D\}\) with \(W' \succeq W\), and for appropriate indices \(i' \in \I_{W'}\), \(i \in \I_W\), there are designated descent morphisms
\[
\iota^{W'W}_{i',i} : X_{i'} \to X_i
\]
in \(\Noetica\), satisfying:
\begin{enumerate}[label=(\roman*)]
\item reflexivity: \(\iota^{WW}_{i,i} = \id_{X_i}\);
\item transitivity: if \(W'' \succeq W' \succeq W\), then whenever the indices are composable,
\[
\iota^{W'W}_{i',i} \circ \iota^{W''W'}_{i'',i'} = \iota^{W''W}_{i'',i}.
\]
\end{enumerate}
\end{axiom}

\begin{axiom}
Let \(i \in \I_W\). For every morphism \(g : J \to X_i\) in \(\Noetica\), the pullbacks
\[
P_{i'} := J \times_{X_i} X_{i'}
\]
along the designated descent maps \(\iota^{W'W}_{i',i} : X_{i'} \to X_i\) with \(W' \succeq W\) exist in \(\Noetica\), and their projection morphisms
\[
q_{i'} : P_{i'} \to J
\]
are jointly surjective on underlying sets:
\[
\bigcup_{i',\,W' \succeq W} \mathrm{Im}\bigl(U(q_{i'})\bigr) = U(J).
\]
\end{axiom}

\begin{remark}
The pullback-stability axiom is the only place where the descent data are required to interact with arbitrary morphisms into a distinguished object. It is precisely the hypothesis needed later to reduce descent coverings to jointly surjective coverings.
\end{remark}

\section{RPM Structure}

\begin{axiom}
There exists an endofunctor
\[
R : \Noetica \to \Noetica.
\]
\end{axiom}

\begin{axiom}
There exist natural transformations
\[
\eta : \id_{\Noetica} \Rightarrow R,
\qquad
\varepsilon : R \Rightarrow \id_{\Noetica}.
\]
\end{axiom}

\begin{axiom}
For every object \(X \in \mathrm{Ob}(\Noetica)\),
\[
\varepsilon_X \circ \eta_X = \id_X.
\]
\end{axiom}

\begin{remark}
The preceding axioms specify a section-retraction pair. No second functor is introduced, and no adjunction or triangle identities are assumed.
\end{remark}

\begin{remark}
Naturality of \(\eta\) and \(\varepsilon\) means that for every morphism \(g : J \to I\) in \(\Noetica\),
\[
R(g)\circ \eta_J = \eta_I \circ g,
\qquad
\varepsilon_I \circ R(g) = g \circ \varepsilon_J.
\]
\end{remark}

\begin{proposition}\label{prop:rpm-pullback}
For any morphism \(g : J \to I\) in \(\Noetica\), let
\[
P := J \times_I R(I)
\]
be the pullback of \(\varepsilon_I : R(I) \to I\) along \(g\), and let \(q : P \to J\) be the projection. Then \(U(q)\) is surjective.
\end{proposition}

\begin{proof}
Define
\[
\phi := \langle \id_J,\, R(g)\circ \eta_J\rangle : J \to P.
\]
To see that \(\phi\) is well defined, we verify the pullback compatibility condition:
\[
\varepsilon_I \circ R(g)\circ \eta_J
= g \circ \varepsilon_J \circ \eta_J
= g \circ \id_J
= g,
\]
where the first equality is naturality of \(\varepsilon\) and the second is Axiom~5.3. Hence \(\phi\) exists by the universal property of the pullback. By construction,
\[
q \circ \phi = \id_J.
\]
Therefore \(q\) has a section. Any function with a section is surjective, so \(U(q)\) is surjective.
\end{proof}

\begin{proposition}\label{prop:finite-model}
The axioms of Sections~3--5 are jointly satisfiable.
\end{proposition}

\begin{proof}
Let
\[
\begin{aligned}
\I&=\{1,2,\dots,40\},\\
\IA&=\{1,\dots,10\},\\
\IB&=\{11,\dots,20\},\\
\IC&=\{21,\dots,30\},\\
\ID&=\{31,\dots,40\}.
\end{aligned}
\]
Let \(1=\{\ast\}\) be a fixed singleton set, and set
\[
X_i:=1
\qquad
\text{for all } i\in \I.
\]
Let
\[
D_0=\{(i,i):i\in \I\},
\qquad
D_k=\I\times \I
\quad
\text{for } 1\leq k\leq 9.
\]
For every admissible pair \((i,j)\in D_k\), define
\[
\nu_{k,i,j}:1\to 1
\]
to be the unique map, namely \(\id_1\).

The category generated by these data is the terminal subcategory of \(\FinSet\) with one object \(1\) and one morphism \(\id_1\). It is concrete, small, and has pullbacks. Indeed, the pullback of \(\id_1\) along \(\id_1\) is again a singleton object, and in \(\FinSet\) this pullback may be represented by \(1\) itself with both projections equal to \(\id_1\). Thus Definition~\ref{def:noetica} and the smallness axiom are satisfied.

For every \(W'\succeq W\), every \(i'\in \I_{W'}\), and every \(i\in \I_W\), define the descent map
\[
\iota^{W'W}_{i',i}:1\to 1
\]
to be \(\id_1\). Reflexivity and transitivity are immediate because all such maps are identities on the unique object. If \(g:J\to X_i\) is a morphism in \(\Noetica\), then \(J=1\) and \(g=\id_1\). The pullbacks required in Section~4 are again represented by \(1\), and each projection \(q_{i'}:1\to 1\) is \(\id_1\), hence surjective. Therefore the ACBE axioms hold.

Finally, let
\[
R:=\id_{\Noetica},
\qquad
\eta:=\id_{\id_{\Noetica}},
\qquad
\varepsilon:=\id_{\id_{\Noetica}}.
\]
Then \(\varepsilon_X\circ \eta_X=\id_X\) for the unique object \(X=1\), and naturality is automatic. Hence the RPM axioms also hold.
\end{proof}

\begin{remark}
In the model of Proposition~\ref{prop:finite-model}, the coverage \(\Jnoe\) is the maximal topology on the terminal category, so the associated sheaf topos is equivalent to \(\mathrm{Set}\). The value of the example is not richness but non-vacuity: the axioms admit an explicit finite realization.
\end{remark}

\section{The Noetic Coverage}

\begin{definition}
For an object \(X \in \mathrm{Ob}(\Noetica)\), a sieve \(S\) on \(X\) is declared covering, written \(S \in \Jnoe(X)\), if at least one of the following holds:
\begin{enumerate}[label=(\textnormal{Noe-Cov-\arabic*})]
\item \label{cov:joint} \(S\) contains a family of morphisms \(\{f_{\alpha} : X_{\alpha} \to X\}_{\alpha \in A}\) whose underlying functions are jointly surjective:
\[
\bigcup_{\alpha \in A}\mathrm{Im}\bigl(U(f_{\alpha})\bigr)=U(X).
\]
\item \label{cov:descent} \(X=X_i\) is a distinguished object for some \(i \in \I_W\), and \(S\) contains all designated descent morphisms into \(X_i\):
\[
\iota^{W'W}_{i',i} : X_{i'} \to X_i
\qquad
(W' \succeq W).
\]
\item \label{cov:rpm} \(S\) contains the counit morphism
\[
\varepsilon_X : R(X) \to X.
\]
\end{enumerate}
\end{definition}

\begin{remark}
Conditions \ref{cov:descent} and \ref{cov:rpm} are typed as genuine sieves on the same codomain \(X\). The proof of Lemma~\ref{lem:reduction} shows that they contribute no additional coverage strength beyond jointly surjective families.
\end{remark}

\begin{lemma}\label{lem:reduction}
If \(S \in \Jnoe(X)\) by condition \ref{cov:descent} or condition \ref{cov:rpm}, then \(S \in \Jnoe(X)\) by condition \ref{cov:joint}.
\end{lemma}

\begin{proof}
Assume first that \(S\) is covering by condition \ref{cov:descent}. Then \(X=X_i\) for some \(i \in \I_W\), and \(S\) contains every designated descent morphism into \(X_i\). Apply Axiom~4.2 with \(g=\id_{X_i}\). The pullbacks
\[
X_i \times_{X_i} X_{i'}
\]
identify canonically with \(X_{i'}\), and the associated projections identify with the descent morphisms \(\iota^{W'W}_{i',i}\). By Axiom~4.2, these projections are jointly surjective on \(U(X_i)\), so \(S\) satisfies condition \ref{cov:joint}.

Now assume that \(S\) is covering by condition \ref{cov:rpm}. Then \(\varepsilon_X \in S\). By Axiom~5.3,
\[
\varepsilon_X \circ \eta_X = \id_X,
\]
so \(U(\varepsilon_X)\) is surjective. Hence the singleton family \(\{\varepsilon_X\}\) is jointly surjective, and \(S\) again satisfies condition \ref{cov:joint}.
\end{proof}

\begin{theorem}\label{thm:Jnoe-topology}
The assignment \(X \mapsto \Jnoe(X)\) is a Grothendieck topology on \(\Noetica\).
\end{theorem}

\begin{proof}
We verify axioms \ref{J1}--\ref{J3}.

\smallskip
\noindent\textbf{(J1) Maximality.}
Let \(X \in \mathrm{Ob}(\Noetica)\). The maximal sieve \(m_X\) contains \(\id_X\), and \(U(\id_X)\) is surjective. Hence \(m_X\) satisfies condition \ref{cov:joint}, so \(m_X \in \Jnoe(X)\).

\smallskip
\noindent\textbf{(J2) Stability under pullback.}
Let \(S \in \Jnoe(X)\) and \(g : J \to X\). By Lemma~\ref{lem:reduction}, it is enough to consider the case in which \(S\) satisfies condition \ref{cov:joint}. Choose morphisms
\[
f_{\alpha} : X_{\alpha} \to X \qquad (\alpha \in A)
\]
in \(S\) such that
\[
\bigcup_{\alpha \in A}\mathrm{Im}\bigl(U(f_{\alpha})\bigr)=U(X).
\]
Because \(\Noetica\) has pullbacks, each diagram
\[
\begin{CD}
P_{\alpha} @>{p_{\alpha}}>> X_{\alpha}\\
@V{q_{\alpha}}VV @VV{f_{\alpha}}V\\
J @>{g}>> X
\end{CD}
\]
exists in \(\Noetica\).

We claim that the family \(\{q_{\alpha} : P_{\alpha} \to J\}_{\alpha \in A}\) is jointly surjective on \(U(J)\). Let \(y \in U(J)\). Then \(U(g)(y)\in U(X)\). By joint surjectivity of the \(f_{\alpha}\), choose \(\alpha \in A\) and \(z \in U(X_{\alpha})\) with
\[
U(f_{\alpha})(z)=U(g)(y).
\]
By the set-theoretic description of pullbacks in \(\FinSet\), the pair \((y,z)\) determines an element of \(U(P_{\alpha})\), and \(U(q_{\alpha})(y,z)=y\). Thus every \(y \in U(J)\) lies in some image \(\mathrm{Im}(U(q_{\alpha}))\), proving the claim.

Each \(q_{\alpha}\) lies in the pullback sieve \(g^{*}S\) because
\[
g \circ q_{\alpha} = f_{\alpha}\circ p_{\alpha},
\]
and \(f_{\alpha}\in S\) implies \(f_{\alpha}\circ p_{\alpha}\in S\) by sieve closure under precomposition. Hence \(g^{*}S\) contains a jointly surjective family and therefore lies in \(\Jnoe(J)\).

\smallskip
\noindent\textbf{(J3) Local character.}
Let \(S \in \Jnoe(X)\), and let \(T\) be a sieve on \(X\) such that for every morphism \(f : D \to X\) in \(S\), the pullback sieve \(f^{*}T\) lies in \(\Jnoe(D)\). We must prove that \(T \in \Jnoe(X)\).

By Lemma~\ref{lem:reduction}, choose a jointly surjective family
\[
f_{\alpha} : X_{\alpha} \to X \qquad (\alpha \in A)
\]
contained in \(S\). For each \(\alpha\), the sieve \(f_{\alpha}^{*}T\) is covering on \(X_{\alpha}\), so by Lemma~\ref{lem:reduction} again it contains a jointly surjective family
\[
h_{\alpha\beta} : K_{\alpha\beta} \to X_{\alpha}
\qquad
(\beta \in B_{\alpha}).
\]
Because \(h_{\alpha\beta} \in f_{\alpha}^{*}T\), every composite
\[
f_{\alpha}\circ h_{\alpha\beta} : K_{\alpha\beta} \to X
\]
lies in \(T\).

We claim that the family \(\{f_{\alpha}\circ h_{\alpha\beta}\}_{\alpha,\beta}\) is jointly surjective on \(U(X)\). Let \(x \in U(X)\). Since the \(f_{\alpha}\) are jointly surjective, choose \(\alpha\) and \(z \in U(X_{\alpha})\) with \(U(f_{\alpha})(z)=x\). Since the \(h_{\alpha\beta}\) are jointly surjective on \(U(X_{\alpha})\), choose \(\beta\) and \(w \in U(K_{\alpha\beta})\) with \(U(h_{\alpha\beta})(w)=z\). Then
\[
U(f_{\alpha}\circ h_{\alpha\beta})(w)
= U(f_{\alpha})(U(h_{\alpha\beta})(w))
= U(f_{\alpha})(z)
= x.
\]
Thus the composites are jointly surjective. Since they lie in \(T\), the sieve \(T\) satisfies condition \ref{cov:joint}. Therefore \(T \in \Jnoe(X)\).
\end{proof}

\section{The Noetic Topos}

\begin{definition}
The \emph{noetic topos} is the category
\[
\NoeTop := \Sh(\Noetica,\Jnoe).
\]
\end{definition}

\begin{remark}
The definition depends only on the small site \((\Noetica,\Jnoe)\). No separate empirical interpretation is needed for the topos-theoretic statement.
\end{remark}

\begin{proposition}
The pair \((\Noetica,\Jnoe)\) is a small site.
\end{proposition}

\begin{proof}
Smallness of \(\Noetica\) was proved in Section~3, and Theorem~\ref{thm:Jnoe-topology} proves that \(\Jnoe\) is a Grothendieck topology.
\end{proof}

\begin{theorem}\label{thm:noetic-topos}
The category \(\NoeTop\) is a Grothendieck topos and therefore an elementary topos. In particular, it has finite limits, finite colimits, exponentials, a subobject classifier \(\Omega\), and intuitionistic internal logic.
\end{theorem}

\begin{proof}
By the preceding proposition, \((\Noetica,\Jnoe)\) is a small site. The general theorem on sheaves over a small site implies that
\[
\Sh(\Noetica,\Jnoe)
\]
is a Grothendieck topos; see \cite[Chap.~VII, Thm.~1]{MacLaneMoerdijk1992}. Every Grothendieck topos is an elementary topos. Consequently \(\NoeTop\) has finite limits, finite colimits, exponentials, and a subobject classifier \(\Omega\). As every elementary topos, it supports an internal higher-order intuitionistic logic \cite[Chap.~VI]{MacLaneMoerdijk1992}.
\end{proof}

\section{Noetic Geometry}

The geometric definitions in this section are standard constructions from information geometry; see \cite{AmariNagaoka2000}.

\begin{definition}
The \emph{open noetic state space} is the open probability simplex
\[
\Deltaopen
:=
\left\{
p=(p_1,\dots,p_{40}) \in \R^{40}
\;\middle|\;
p_k>0 \text{ for all } k,\;
\sum_{k=1}^{40} p_k = 1
\right\}.
\]
It is a smooth manifold of dimension \(39\).
\end{definition}

\begin{definition}
A \emph{noetic manifold} is a pair \((M,\sigma)\) consisting of a smooth manifold \(M\) and a smooth map
\[
\sigma : M \to \Deltaopen.
\]
\end{definition}

\begin{definition}
For a noetic manifold \((M,\sigma)\) and a point \(p \in M\), the \emph{noetic dimension} at \(p\) is
\[
\dim_{\sigma}(p) := \dim(T_pM) + \mathrm{rank}(d\sigma_p).
\]
This is a bookkeeping quantity attached to the pair \((M,\sigma)\); it is not claimed to be a standard invariant of differential geometry.
\end{definition}

\begin{definition}
The \emph{Fisher--Rao metric} on \(\Deltaopen\) is the Riemannian metric
\[
g^{\mathrm{FR}}_p(u,v)
:=
\sum_{k=1}^{40}\frac{u_kv_k}{p_k},
\]
for \(p \in \Deltaopen\) and tangent vectors \(u,v \in T_p\Deltaopen = \{x \in \R^{40} : \sum_k x_k = 0\}\).
The pullback bilinear form on a noetic manifold \((M,\sigma)\) is
\[
g^{\sigma}_p(\xi,\zeta)
:=
g^{\mathrm{FR}}_{\sigma(p)}\bigl(d\sigma_p(\xi),d\sigma_p(\zeta)\bigr).
\]
\end{definition}

\begin{example}
Taking \(M=\Deltaopen\) and \(\sigma=\id_{\Deltaopen}\) yields a canonical noetic manifold. Since \(\dim(\Deltaopen)=39\) and \(d(\id_{\Deltaopen})_p\) has rank \(39\) at every point, the noetic dimension in this example is \(78\).
\end{example}

\begin{remark}
The geometric layer is independent of the categorical proofs in Sections~3--7. Its role is to provide a smooth state-space language that can be paired with the site/topos construction if one later introduces empirical probability assignments.
\end{remark}

\section{Discussion}

The paper proves a conditional mathematical statement. If the finite distinguished data, ACBE descent axioms, and RPM section-retraction axioms are adopted, then the induced coverage \(\Jnoe\) is a Grothendieck topology and the corresponding sheaf category is a Grothendieck topos. Proposition~\ref{prop:finite-model} shows that these axioms are not empty formal placeholders: they admit at least one explicit finite realization. The Fisher--Rao construction then supplies a compatible information-geometric layer on the open simplex. That is the full scope of the present result.

Equally important is what the paper does \emph{not} prove. It does not show that the axioms are forced by first principles, that they are uniquely natural, or that they have yet been validated against empirical measurements of AI behavior. The axioms in Sections~3--5 are proposed formal inputs. Any claim that they support a real assessment methodology requires future empirical work: dataset construction, operationalization of observables, calibration, reliability analysis, predictive validation, and adversarial testing.

With that caveat, the framework has a plausible future use. A structured assessment program for AI systems often collects partially overlapping local observations: benchmark outcomes, prompt-conditioned safety failures, interpretability measurements, and deployment-context annotations. The site/topos language developed here could, in principle, serve as a disciplined way to index such local views, specify which families count as jointly informative, and distinguish compatible from incompatible gluing.

This positioning is compatible with recent empirical interpretability work rather than a substitute for it. Anthropic's 2026 study of functional emotion concepts in Claude reports behaviorally relevant internal representations and explicitly avoids claims about subjective experience \cite{Anthropic2026Emotion}. The present paper does not explain those findings and does not compete with them. At most, it suggests a future formal layer for organizing multi-context evidence of that general kind, should one wish to relate local interpretability signals to broader assessment structures.

\bibliographystyle{plain}
\bibliography{references}

\end{document}
